%% Capítulo 1: Introducción

\chapter{Introducción}
\label{chap:introduccion}

Los modelos de equilibrio general dinámico estocástico (DSGE, por sus siglas en inglés)
se convirtieron en el estándar de análisis macroeconómico de los bancos centrales durante
las últimas décadas \citep{SmetsWouters2007}. Estos modelos basan su poder teórico en sus bases
microfundadas de las cuales los comportamientos agregados se derivan a partir de la optimimzación
de los diferentes agentes económicos representativos (hogares y firmas) bajo expectativas
racionales, lo que permite una interpretación estructural
de los parámetros y, en principio, la evaluación de políticas económicas sin caer en la
crítica de Lucas \citep{Lucas1976}. De esta forma los modelos DSGE, aislan los parámetros estructurales
como ser las preferencias de consumo o la tecnología, que se asumen estables aunque cambie la políca del gobierno.
Sin embargo, la crísis del 2008-2009 expuso las limitaciones de
este paradigma con particular crudeza: los modelos DSGE estándar fallaron en anticipar
la crisis y en replicar su dinámica, al tiempo que las fricciones financieras y la
heterogeneidad de los agentes demostraron ser determinantes, no ornamentales
\citep{FagioloRoventini2017}.



\section{Motivación}
\label{sec:motivacion}

\section{Objetivos}
\label{sec:objetivos}

\section{Contribuciones}
\label{sec:contribuciones}

\section{Estructura del documento}
\label{sec:estructura}
El presente documento cuenta con tres secciones: una sección general, donde se incluyen 
introduccion, Marco teórico y estado del arte. Una sección donde se plantea una Investigación
académica y finalmente una última sección donde se extrapola la investicación académica como un 
producto de negocios.
El Capítulo~\ref{chap:marco-teorico} presenta los tres pilares teóricos que sustentan la
comparación: los fundamentos de los modelos DSGE, el paradigma de los modelos basados en
agentes (ABM) como respuesta a las limitaciones del DSGE, y el uso de LLMs como motores
de decisión para agentes económicos. El Capítulo~\ref{chap:estado-del-arte} revisa
sistemáticamente la literatura en cada uno de esos frentes, con énfasis en los trabajos
sobre Argentina, y cierra identificando los cinco gaps que esta tesis aborda.

El Capítulo~\ref{chap:metodologia} describe el diseño empírico: las fuentes de datos y
el pipeline de acceso, la implementación del modelo DSGE de referencia en Python, y la
arquitectura del sistema de agentes LLM. Los Capítulos~\ref{chap:resultados-dsge}
y~\ref{chap:resultados-agentes} presentan los resultados de cada enfoque por separado.
El Capítulo~\ref{chap:comparacion} integra ambos mediante el protocolo de benchmark y
analiza la respuesta de cada sistema ante los shocks propuestos. El
Capítulo~\ref{chap:conclusiones} sintetiza los hallazgos, evalúa las contribuciones
frente a los gaps identificados, y propone líneas de trabajo futuro.
